\subsection[K-theory of the Roe algebra]{$\boldsymbol{\K}$-theory of the Roe algebra}\label{sec:RoeAlg}

The Roe algebra is a \textCstar-algebra $\Roe(X)$ that can be constructed for every coarse space $X$ (cf.\ \cite[Section~4.4]{RoeCoarseGeometry}), but it is usually only dealt with in the case of proper metric spaces. As most of the references only consider the latter special case, we shall stick to it.

Let us briefly recall the broad outlines of its construction as presented, e.g., in \cite[Section~6.3]{HigRoe}. Starting with a representation $\rho$ of $\Cz(X)$ on a separable Hilbert space $H$ which is non-degenerate (meaning that $\rho(\Cz(X))H$ is dense in $H$), the Roe algebra $\Roe(X,\rho)$ is defined as the norm closure of the $*$-subalgebra of $\Lin(H)$ consisting of all operators $T$ satisfying the following two properties:\looseness=1
\begin{enumerate}\itemsep=0pt
\item[$\bullet$] $T$ is called \emph{locally compact}, if $\rho(f)T,T\rho(f)$ are compact operators on $H$ for all $f\in\Cz(X)$.
\item[$\bullet$] $T$ has \emph{finite propagation}, if there is $R\geq 0$ such that $\rho(f)T\rho(g)=0$ whenever the distance between the supports of $f,g\in\Cz(X)$ is at least $R$. The infimum of all $R\geq 0$ with this property is called the \emph{propagation} $\Prop(T)$ of $T$.
\end{enumerate}
This \textCstar-algebra depends on the choice of $H$ and $\rho$, but if the representation is big enough in a~cer\-tain sense (it is \enquote{ample}, i.e., it is non-degenerate and $\rho^{-1}(\Kom(H))=\{0\}$), then its $\K$-theory $\K_*(\Roe(X,\rho))$ is independent from the choices up to canonical isomorphism.
Therefore, the representation $\rho$ is usually omitted from the notation with the understanding that one suitable representation has been chosen for each $X$.

Now, there are two well-known generalizations of the Roe algebra. The first is the equivariant Roe algebra in the case that a countable discrete group $\Gamma$ acts properly and isometrically on the coarsely proper metric space $X$. Here one chooses in addition to the non-degenerate representation~$\rho$ also a representation $u$ of the group $\Gamma$ on $H$ by unitaries such that $u(\gamma)\rho(f)u(\gamma)^*=\rho(\gamma\cdot f)$ for all $\gamma\in\Gamma$ and $f\in\Cz(X)$. The triple $(H,\rho,u)$ is called a $\Gamma$-$X$-module. Then the equivariant Roe algebra $\Roe_\Gamma(X,\rho,u)$ is the norm closure of the \textCstar-algebra of all locally compact operators of finite propagation in $\Lin(H)$ that are equivariant with respect to $u$. Again, the $\K$-theory groups $\K_*(\Roe_\Gamma(X,\rho,u))$ are independent up to canonical isomorphism from the $\Gamma$-$X$-module $(H,\rho,u)$ if the latter has been chosen sufficiently large and then abbreviated by $\K_*(\Roe_\Gamma(X))$.
A detailed exposition about equivariant Roe algebras can be found in \cite{WillettYuHigherIndexTheory}.

The second generalization is to introduce a (possibly graded) coefficient \textCstar-algebra $D$ into the Roe algebra, cf.\ \cite{HankePapeSchick,HigsonPedersenRoe,WulffTwisted}. This is done by replacing the Hilbert space $H$ by a right Hilbert module $\HilbertMod$ over $D$ and the non-degenerate representation $\rho$ is by adjointable operators on $\HilbertMod$.
The $\K$-theory of the resulting \textCstar-algebra $\Roe(X,\rho,\HilbertMod;D)$ is again unique up to canonical isomorphism if $\rho$ and $\HilbertMod$ have been chosen sufficiently large and in this case $\rho$ will be ommited from the notation. Here, sufficient largeness means that $\HilbertMod\cong H\otimes D$ and $\rho=\rho'\otimes \id_D$ for an ample representation~$\rho'$ on a separable Hilbert space $H$ \cite[Definition 4.5]{HigsonPedersenRoe}.


It is possible to unify the two generalizations to equivariant Roe algebras $\Roe_\Gamma(X,\rho,u,\HilbertMod;D)$ with coefficients in a (graded) $\Gamma$-\textCstar-algebra $D$.
In this case, $u$ also has to be a representation by adjointable operators on $\HilbertMod$ which is compatible with the $\Gamma$-action on $D$ in the sense that $u(\gamma)(\xi d)=(u(\gamma)\xi) (\gamma d)$ for all $\xi\in\HilbertMod$, $d\in D$ and $\gamma\in\Gamma$.
Again, we obtain unambiguous $\K$-theory groups $\K_*(\Roe_\Gamma(X;D))$ if we start with sufficiently large $H,\rho,u$ as in the equivariant case and then simply use their tensor product with $D$, i.e., $\HilbertMod\coloneqq H\otimes D$ with obvious induced representations of $\Cz(X)$ and $\Gamma$.




The proofs of the following statements can be found in the above-mentioned sources in the special cases just described, but the proofs generalize to the general case of equivariant Roe-algebras with coefficients.

If $A\subset X$ is a $\Gamma$-invariant closed subspace of the proper metric $\Gamma$-space $X$, then $\Roe_\Gamma(X;D)$ contains an ideal $\Roe_\Gamma(A\subset X;D)$ whose $\K$-theory is canonically isomorphic to the $\K$-theory of $\Roe_\Gamma(A;D)$.
It is the norm closure of all $T\in\Roe_\Gamma(X;D)$ which are supported in an $R$-neighborhood of $A$ for some $R\geq 0$, that is, all those $T$ for which $\rho(f)T=T\rho(f)=0$ for all $f\in\Cz(X)$ with $\dist(\supp(f),A)\geq R$.
The quotient \textCstar-algebra $\Roe_\Gamma(X,A;D)\coloneqq \Roe_\Gamma(X;D)/\Roe_\Gamma(A\subset X;D)$ is called the \emph{$\Gamma$-equivariant relative Roe algebra}. We have $\Roe_\Gamma(\varnothing\subset X;D)=0$ and hence $\Roe_\Gamma(X,\varnothing;D)=\Roe_\Gamma(X;D)$. The long exact sequences
\[
\dots\to\K_p(\Roe_\Gamma(A;D))\to \K_p(\Roe_\Gamma(X;D))\to \K_p(\Roe_\Gamma(X,A;D))\to \K_{p-1}(\Roe_\Gamma(A;D))\to \cdots
\]
are immediate.
